% GENERATED FILE. Edit canonical lessons, then run npm run generate:books.
% canonical-source-hash: fnv1a64:9de86ae13ca666cb
% canonical-lessons: PA-C195-window-room-return, PA-C195-vala, PA-C195-vala-review, PA-C195-vale, PA-C195-vale-review, PA-C195-direct-oblique-contrast, PA-C195-payoff

\chapter{A Room with a Window}
\label{ch:pa-a2-vala-forms}
\hlchaptermodality{195}

\begin{chapteropening}
\textbf{By the end of this chapter:} \emph{I can contrast a room with a window with in that room, using one bounded masculine direct-oblique pattern.}

Contrast a room with a window and in that room through listening, speaking, reading, and writing.
\end{chapteropening}

\section[\texorpdfstring{\pa{ਖਿੜਕੀ} · \pa{ਕਮਰਾ} … (khirkii; kamraa …)}{ਖਿੜਕੀ · ਕਮਰਾ … (khirkii; kamraa …)}]{Return to window and room}

\label{lesson:PA-C195-window-room-return}



\begin{quote}
\emph{Hear:} \emph{khirkii; kamraa}; YOU SAY: \textbf{\pa{ਖਿੜਕੀ}; \pa{ਕਮਰਾ}}
\end{quote}

\begin{grammarlens}[title={Grammar Lens}]
Keep the familiar room contrast small: \textbf{\pa{ਕਮਰਾ}} on its own; \textbf{\pa{ਕਮਰੇ} \pa{ਵਿੱਚ}} for “in the room.” This opener adds no new form.
\end{grammarlens}

\subsection*{Your turn}
\begin{itemize}
\raggedright
  \item \emph{Hear:} \emph{khirkii}; YOU SAY: \textbf{\pa{ਖਿੜਕੀ}}
  \item \emph{Hear:} \emph{kamraa}; YOU SAY: \textbf{\pa{ਕਮਰਾ}}
  \item \emph{Hear:} \emph{kamre vich}; YOU SAY: \textbf{\pa{ਕਮਰੇ} \pa{ਵਿੱਚ}}
\end{itemize}

\begin{tcolorbox}[breakable,colback=teal!4,colframe=teal!35!black,title={Before you move on}]
Say “window,” “room,” and “in the room.”
\end{tcolorbox}
\section[\texorpdfstring{\pa{ਖਿੜਕੀ} \pa{ਵਾਲਾ} \pa{ਕਮਰਾ} (khirkii vaalaa kamraa)}{ਖਿੜਕੀ ਵਾਲਾ ਕਮਰਾ (khirkii vaalaa kamraa)}]{Add one bounded association to room}

\label{lesson:PA-C195-vala}



\begin{quote}
Retrieve \textbf{\pa{ਖਿੜਕੀ}}, then \textbf{\pa{ਕਮਰਾ}}.
\end{quote}

\begin{grammarlens}[title={Grammar Lens: direct owner}]
\begin{quote}
\textbf{\pa{ਖਿੜਕੀ} \pa{ਵਾਲਾ} \pa{ਕਮਰਾ}} — \emph{khiṛkī vālā kamrā} — a room with a window
\end{quote}

Here noun \textbf{\pa{ਖਿੜਕੀ}} plus \textbf{\pa{ਵਾਲਾ}} forms an adjective phrase describing the masculine singular noun \textbf{\pa{ਕਮਰਾ}}. Learn this one direct room phrase before any broader pattern.
\end{grammarlens}

\begin{grammarlens}[title={Grammar Lens: source and scope boundary}]
Gill and Gleason describe noun plus \textbf{\pa{ਵਾਲਾ}} as an adjective phrase and show the masculine adjective direct–oblique pattern. This chapter follows only one room contrast. It is not a rule for every noun or postposition. The earlier \textbf{\pa{ਵਾਲੀ}}, “earring,” word lesson owns a separate lexical homograph, not a grammatical form licensed by that old word lesson. Source: Harjeet Singh Gill and Henry A. Gleason, Jr., \emph{A Reference Grammar of Punjabi} (1969), pp. 71–72 and 107, \href{https://pt.learnpunjabi.org/assets/A\%20Reference\%20Grammar\_Final.pdf}{adjective phrases}.
\end{grammarlens}

\subsection*{Your turn}
\begin{itemize}
\raggedright
  \item \emph{Hear:} \emph{khirkii vaalaa kamraa}; YOU SAY: \textbf{\pa{ਖਿੜਕੀ} \pa{ਵਾਲਾ} \pa{ਕਮਰਾ}}
  \item \emph{Read it:} \textbf{\pa{ਖਿੜਕੀ} \pa{ਵਾਲਾ} \pa{ਕਮਰਾ}}
  \item \emph{Write it:} copy \textbf{\pa{ਵਾਲਾ}} once, then rebuild the whole phrase
\end{itemize}

\begin{tcolorbox}[breakable,colback=teal!4,colframe=teal!35!black,title={Before you move on}]
Produce “a room with a window.”
\end{tcolorbox}
\section[\texorpdfstring{\pa{ਖਿੜਕੀ} \pa{ਵਾਲਾ} \pa{ਕਮਰਾ} (khirkii vaalaa kamraa)}{ਖਿੜਕੀ ਵਾਲਾ ਕਮਰਾ (khirkii vaalaa kamraa)}]{Retrieve the direct room phrase}

\label{lesson:PA-C195-vala-review}



\begin{quote}
Say the new direct form \textbf{\pa{ਵਾਲਾ}}.
\end{quote}

\subsection*{Your turn}
\begin{itemize}
\raggedright
  \item \emph{Hear:} \emph{khirkii vaalaa kamraa}; YOU SAY: \textbf{\pa{ਖਿੜਕੀ} \pa{ਵਾਲਾ} \pa{ਕਮਰਾ}}
  \item \emph{Read it:} \textbf{\pa{ਖਿੜਕੀ} \pa{ਵਾਲਾ} \pa{ਕਮਰਾ}}
  \item \emph{Write it:} cover and write \textbf{\pa{ਵਾਲਾ}}
\end{itemize}

\begin{grammarlens}[title={Grammar Lens}]
Keep the meaning and form together inside this one room phrase.
\end{grammarlens}

\begin{tcolorbox}[breakable,colback=teal!4,colframe=teal!35!black,title={Before you move on}]
Produce “a room with a window.”
\end{tcolorbox}
\section[\texorpdfstring{\pa{ਖਿੜਕੀ} \pa{ਵਾਲੇ} \pa{ਕਮਰੇ} \pa{ਵਿੱਚ} (khirkii vaale …)}{ਖਿੜਕੀ ਵਾਲੇ ਕਮਰੇ ਵਿੱਚ (khirkii vaale …)}]{Change the association inside one location frame}

\label{lesson:PA-C195-vale}



\begin{quote}
Retrieve \textbf{\pa{ਖਿੜਕੀ} \pa{ਵਾਲਾ} \pa{ਕਮਰਾ}}, then \textbf{\pa{ਕਮਰੇ} \pa{ਵਿੱਚ}}.
\end{quote}

\begin{grammarlens}[title={Grammar Lens: direct owner}]
\begin{quote}
\textbf{\pa{ਖਿੜਕੀ} \pa{ਵਾਲਾ} \pa{ਕਮਰਾ}} — \emph{khiṛkī vālā kamrā} — a room with a window
\end{quote}

\begin{quote}
\textbf{\pa{ਖਿੜਕੀ} \pa{ਵਾਲੇ} \pa{ਕਮਰੇ} \pa{ਵਿੱਚ}} — \emph{khiṛkī vāle kamre vicc} — in the room with a window
\end{quote}

In this second phrase, \textbf{\pa{ਵਾਲੇ}} is masculine singular oblique. It agrees with masculine singular oblique \textbf{\pa{ਕਮਰੇ}} before \textbf{\pa{ਵਿੱਚ}}. Here \textbf{\pa{ਵਾਲੇ}} is not a plural or honorific form. This chapter does not extend the pattern to every noun or postposition.
\end{grammarlens}

\subsection*{Your turn}
\begin{itemize}
\raggedright
  \item \emph{Hear:} \emph{khirkii vaale kamre vich}; YOU SAY: \textbf{\pa{ਖਿੜਕੀ} \pa{ਵਾਲੇ} \pa{ਕਮਰੇ} \pa{ਵਿੱਚ}}
  \item \emph{Read it:} \textbf{\pa{ਖਿੜਕੀ} \pa{ਵਾਲੇ} \pa{ਕਮਰੇ} \pa{ਵਿੱਚ}}
  \item \emph{Write it:} copy \textbf{\pa{ਵਾਲੇ}} once, then compare it with \textbf{\pa{ਵਾਲਾ}}
\end{itemize}

\begin{tcolorbox}[breakable,colback=teal!4,colframe=teal!35!black,title={Before you move on}]
Produce “in the room with a window.”
\end{tcolorbox}
\section[\texorpdfstring{\pa{ਖਿੜਕੀ} \pa{ਵਾਲੇ} \pa{ਕਮਰੇ} \pa{ਵਿੱਚ} (khirkii vaale …)}{ਖਿੜਕੀ ਵਾਲੇ ਕਮਰੇ ਵਿੱਚ (khirkii vaale …)}]{Retrieve the location phrase}

\label{lesson:PA-C195-vale-review}



\begin{quote}
Say the new form \textbf{\pa{ਵਾਲੇ}}.
\end{quote}

\subsection*{Your turn}
\begin{itemize}
\raggedright
  \item \emph{Hear:} \emph{khirkii vaale kamre vich}; YOU SAY: \textbf{\pa{ਖਿੜਕੀ} \pa{ਵਾਲੇ} \pa{ਕਮਰੇ} \pa{ਵਿੱਚ}}
  \item \emph{Read it:} \textbf{\pa{ਖਿੜਕੀ} \pa{ਵਾਲੇ} \pa{ਕਮਰੇ} \pa{ਵਿੱਚ}}
  \item \emph{Write it:} cover and write \textbf{\pa{ਵਾਲੇ}}
\end{itemize}

\begin{grammarlens}[title={Grammar Lens}]
Keep \textbf{\pa{ਵਾਲੇ}} beside \textbf{\pa{ਕਮਰੇ} \pa{ਵਿੱਚ}} in this bounded masculine singular oblique location frame.
\end{grammarlens}

\begin{tcolorbox}[breakable,colback=teal!4,colframe=teal!35!black,title={Before you move on}]
Produce the complete location phrase.
\end{tcolorbox}
\section[\texorpdfstring{\pa{ਖਿੜਕੀ} \pa{ਵਾਲਾ} \pa{ਕਮਰਾ} … (khirkii vaalaa …)}{ਖਿੜਕੀ ਵਾਲਾ ਕਮਰਾ … (khirkii vaalaa …)}]{Put the two room phrases side by side}

\label{lesson:PA-C195-direct-oblique-contrast}



\begin{quote}
Retrieve \textbf{\pa{ਵਾਲਾ}}, then \textbf{\pa{ਵਾਲੇ}}.
\end{quote}

\begin{grammarlens}[title={Grammar Lens}]
\noindent
\begin{tabularx}{\linewidth}{@{}>{\raggedright\arraybackslash}X>{\raggedright\arraybackslash}X@{}}
\toprule
\textbf{Direct room phrase} & \textbf{Oblique location phrase} \\
\midrule
\textbf{\pa{ਖਿੜਕੀ} \pa{ਵਾਲਾ} \pa{ਕਮਰਾ}} & \textbf{\pa{ਖਿੜਕੀ} \pa{ਵਾਲੇ} \pa{ਕਮਰੇ} \pa{ਵਿੱਚ}} \\
\bottomrule
\end{tabularx}

Read the direct phrase first. Then move to the one familiar location frame. The chapter claims only this contrast.
\end{grammarlens}

\subsection*{Your turn}
\begin{itemize}
\raggedright
  \item \emph{Hear:} \emph{khirkii vaalaa kamraa}; YOU SAY: \textbf{\pa{ਖਿੜਕੀ} \pa{ਵਾਲਾ} \pa{ਕਮਰਾ}}
  \item \emph{Hear:} \emph{khirkii vaale kamre vich}; YOU SAY: \textbf{\pa{ਖਿੜਕੀ} \pa{ਵਾਲੇ} \pa{ਕਮਰੇ} \pa{ਵਿੱਚ}}
  \item \emph{Write it:} cover the table and write both complete phrases
\end{itemize}

\begin{tcolorbox}[breakable,colback=teal!4,colframe=teal!35!black,title={Before you move on}]
Say the direct phrase, then the location phrase.
\end{tcolorbox}
\section[\texorpdfstring{\pa{ਖਿੜਕੀ} \pa{ਵਾਲਾ} \pa{ਕਮਰਾ} … (khirkii vaalaa …)}{ਖਿੜਕੀ ਵਾਲਾ ਕਮਰਾ … (khirkii vaalaa …)}]{A bounded room-association payoff}

\label{lesson:PA-C195-payoff}



\begin{quote}
Recall the direct form, then the form in the location frame.
\end{quote}

\subsection*{Your turn: four skills}
\begin{itemize}
\raggedright
  \item \emph{Hear:} \emph{khirkii vaalaa kamraa; khirkii vaale kamre vich}; YOU SAY: \textbf{\pa{ਖਿੜਕੀ} \pa{ਵਾਲਾ} \pa{ਕਮਰਾ}; \pa{ਖਿੜਕੀ} \pa{ਵਾਲੇ} \pa{ਕਮਰੇ} \pa{ਵਿੱਚ}}
  \item \emph{Read it:} \textbf{\pa{ਖਿੜਕੀ} \pa{ਵਾਲਾ} \pa{ਕਮਰਾ} · \pa{ਖਿੜਕੀ} \pa{ਵਾਲੇ} \pa{ਕਮਰੇ} \pa{ਵਿੱਚ}}
  \item \emph{Write it:} cover the line and write both phrases
\end{itemize}

\begin{tcolorbox}[breakable,colback=teal!4,colframe=teal!35!black,title={Before you move on}]
Produce both complete room phrases.

This is the eighteenth prerequisite tranche exposed by Punjabi A2 mock drafting. It does not unblock Punjabi A2 mock 1 by itself. It is not an A2 mock, not A2 readiness or attainment, and not external or human validation.
\end{tcolorbox}
